\documentclass[11pt,a4paper,fontset=windows]{ctexart}
% ===== 公共导言区：高等代数（一）期中模拟卷 =====
\usepackage{amsmath,amssymb}
\usepackage{unicode-math}
\setmathfont{TeX Gyre Termes Math}
\usepackage[margin=2.05cm,top=2.35cm,bottom=2.0cm,headsep=11pt,includefoot]{geometry}
\usepackage{fancyhdr}
\usepackage{lastpage}
\usepackage{enumitem}
\usepackage{array}
\usepackage{bm}
\usepackage{xcolor}
\usepackage{zhnumber}
\usepackage{needspace}

\definecolor{hdgray}{gray}{0.32}
\definecolor{lngray}{gray}{0.5}
\definecolor{notebg}{gray}{0.955}

\pagestyle{fancy}
\fancyhf{}
\renewcommand{\headrulewidth}{0.4pt}
\renewcommand{\footrulewidth}{0pt}
\fancyhead[L]{\small\color{hdgray}\songti 高等代数（一）\ 期中模拟试卷（\thepp）}
\fancyhead[R]{\small\color{hdgray}\songti 考试时间：120 分钟\quad 满分：100 分}
\fancyfoot[C]{\small\color{lngray}\songti 第 \thepage\ 页\quad 共 \pageref{LastPage}\ 页}

\setlength{\parindent}{0pt}
\setlength{\parskip}{2pt}
\linespread{1.16}

% 题目列表 1—8 连续编号
\newlist{qs}{enumerate}{1}
\setlist[qs]{label=\textbf{\arabic*.},ref=\arabic*,leftmargin=2.55em,
  labelsep=0.55em,itemsep=23pt,topsep=6pt,parsep=0pt}
\newlist{note}{enumerate}{1}
\setlist[note]{label=(\arabic*),leftmargin=1.75em,itemsep=1.5pt,topsep=2.5pt,parsep=0pt}
\newlist{sub}{enumerate}{1}
\setlist[sub]{label=(\arabic*),leftmargin=1.9em,itemsep=3pt,topsep=3pt,parsep=0pt}

\newcommand{\secA}[1]{\Needspace{9\baselineskip}\par\vspace{10pt}{\large\bfseries\songti #1}\par\vspace{3pt}\nopagebreak}
\newcommand{\blank}[1][2.6cm]{\underline{\hspace{#1}}}
\newcounter{pp}
\renewcommand{\thepp}{\chinese{pp}}

\setcounter{pp}{6}
\begin{document}

\begin{center}
{\LARGE\bfseries\songti 高等代数（一）期中模拟试卷（六）}\\[7pt]
{\normalsize\songti 适用范围：第一章\ 多项式（\S1\,--\,\S9）；第二章\ 行列式（\S2.1\,--\,\S2.7）}\\[3pt]
{\small\color{hdgray}\songti 多元多项式、对称多项式不在考查范围}
\end{center}
\vspace{-4pt}

\begin{center}
\begin{tabular}{@{}p{0.335\textwidth}p{0.335\textwidth}p{0.26\textwidth}@{}}
\songti 姓名：\blank[2.5cm] & \songti 学号：\blank[2.5cm] & \songti 得分：\blank[1.9cm]
\end{tabular}
\end{center}

\vspace{2pt}
\hrule height 0.5pt
\begin{center}
\setlength{\fboxsep}{6pt}%
\colorbox{notebg}{\parbox{0.94\textwidth}{\small\songti
\textbf{考生须知：}
\begin{note}
\item 本试卷共 8 题，满分 100 分，考试时间 120 分钟。
\item 计算题须写出必要的计算过程，证明题须给出完整的推理步骤；只写答案不给分。
\item 所有答案一律写在答题纸上，写在试卷上无效。
\end{note}}}
\end{center}
\hrule height 0.5pt
\vspace{4pt}

\secA{一、计算题（每小题 10 分，共 40 分）}
\begin{qs}
\item 计算 $4$ 阶行列式，并将结果分解为一次因式之积：
\[
D=\begin{vmatrix}
x&a&b&c\\
a&x&c&b\\
b&c&x&a\\
c&b&a&x
\end{vmatrix}.
\]

\item 计算 $n$ 阶行列式
\[
D_n=\begin{vmatrix}
0&1&0&\cdots&0\\
1&0&1&\cdots&0\\
0&1&0&\ddots&\vdots\\
\vdots&\ddots&\ddots&\ddots&1\\
0&\cdots&0&1&0
\end{vmatrix},
\]
并就 $n$ 为奇数、$n$ 为偶数两种情形分别写出结果。

\item 设 $f(x)=x^{6}-1$。
\begin{sub}
\item 写出 $f(x)$ 在 $\mathbb{Q}[x]$、$\mathbb{R}[x]$、$\mathbb{C}[x]$ 中的标准分解式；
\item 判断 $f(x)$ 有无重因式，并说明理由。
\end{sub}

\item 设 $f(x)=2x^{4}-7x^{3}+4x^{2}+7x-6$。
\begin{sub}
\item 求 $f(x)$ 在 $\mathbb{Q}[x]$ 中的标准分解式；
\item 判断 $f(x)$ 有无重因式。
\end{sub}
\end{qs}

\secA{二、证明题（每小题 15 分，共 60 分）}
\begin{qs}[resume*]
\item 设
\[
f(x)=x^{n}+a_1x^{n-1}+\cdots+a_{n-1}x+a_n\in\mathbb{Z}[x]
\]
是首一整系数多项式，$p$ 为素数。
\begin{sub}
\item 证明：若将 $f(x)$ 的各系数模 $p$ 后所得的多项式不能分解为两个次数较低的多项式之积，
      则 $f(x)$ 在 $\mathbb{Q}[x]$ 上不可约；
\item 举例说明：上述命题中“首一”的条件不能去掉。
\end{sub}

\item 设 $f(x)\in P[x]$，$\deg f=n\ge 1$。证明：
\[
f(x)\ \text{没有重因式}\iff \bigl(f(x),\,f'(x)\bigr)=1 .
\]

\item 设 $f(x)\in\mathbb{Q}[x]$ 在 $\mathbb{Q}$ 上不可约，$\deg f\ge 1$，$\alpha\in\mathbb{C}$ 是 $f(x)$ 的一个根
（即 $f(\alpha)=0$）。证明：若 $g(x)\in\mathbb{Q}[x]$ 满足 $g(\alpha)=0$，则
\[
f(x)\mid g(x).
\]

\item 设 $A=(a_{ij})$ 为 $n$ 阶实方阵，满足 $a_{ji}=-a_{ij}$（$1\le i,j\le n$），即 $A$ 是反对称矩阵。
\begin{sub}
\item 证明：当 $n$ 为奇数时，$|A|=0$；
\item 当 $n=4$ 时，证明
\[
|A|=\bigl(a_{12}a_{34}-a_{13}a_{24}+a_{14}a_{23}\bigr)^{2}.
\]
\end{sub}
\end{qs}

\end{document}
